\documentclass[11pt]{article}

% ---- packages ----
\usepackage[letterpaper,margin=1in]{geometry}
\usepackage{microtype}
\usepackage{amsmath}
\usepackage{booktabs}
\usepackage{array}
\usepackage{graphicx}
\usepackage{xcolor}
\usepackage[most]{tcolorbox}
\usepackage{enumitem}
\usepackage{caption}
\usepackage[hidelinks]{hyperref}
\usepackage{parskip}
\usepackage{tocloft}
\setlength{\cftbeforesecskip}{1pt}
\usepackage[nameinlink]{cleveref}

\graphicspath{{figures/}}
\captionsetup{font=small,labelfont=bf}
\setlist{leftmargin=1.4em,topsep=3pt,itemsep=2pt}

% ---- teaching / takeaway callout boxes ----
\definecolor{bgblue}{HTML}{2563EB}
\definecolor{bggreen}{HTML}{15803D}
\newtcolorbox{background}[1]{
  breakable, colback=bgblue!5, colframe=bgblue!60, coltitle=bgblue!75!black,
  fonttitle=\bfseries\small, title={Background you'll need --- #1},
  boxrule=0.6pt, left=6pt, right=6pt, top=4pt, bottom=4pt, arc=2pt}
\newtcolorbox{takeaway}[1]{
  breakable, colback=bggreen!6, colframe=bggreen!55, coltitle=bggreen!60!black,
  fonttitle=\bfseries\small, title={#1},
  boxrule=0.6pt, left=6pt, right=6pt, top=4pt, bottom=4pt, arc=2pt}

\newcommand{\code}[1]{\texttt{\small #1}}

\title{\textbf{The Benchmark Chooses the Winner}\\[2pt]
\large An Annotated, Plain-Language Edition\\[2pt]
\normalsize\itshape Measuring what fine-tuning does to small safety guards --- explained
for readers new to the metrics}
\author{Reza Rahimi \\ \small JazzX AI, Los Altos, CA \\ \small\texttt{reza.rahimi@jazzx.ai}}
\date{}

\begin{document}
\maketitle

\begin{abstract}
\noindent
Teams often take a small chat model, do a quick fine-tune to turn it into a
\emph{safety guard} (a filter that flags harmful prompts), and report how well it scores
on a public benchmark. This edition re-tells a study that asks a sharper question:
\emph{what did the fine-tune change compared to the exact same model before fine-tuning,
and does that change hold up on datasets the model never trained on?} On four small
instruction models (1.5--4B parameters), each fine-tuned five times with LoRA, we find a
clear split. On test data from the \textbf{same datasets used in training} (``represented''
sources), fine-tuning is a large, reliable win: the main score rises by about
\textbf{$+0.33$} (to $\approx0.98$ for every model). On \textbf{four brand-new datasets}
(``transfer''), the same fine-tuning is a small average loss (about \textbf{$-0.05$}), and
at a deployed threshold it catches about \emph{one in five fewer} attacks on a hard
held-out harm set (HarmBench recall $78.4\%\!\to\!57.5\%$). Crucially, the average hides a
pattern: fine-tuning \emph{helps weak base models and hurts strong ones} on new data. The
lesson --- and the title --- is that \textbf{whichever benchmark you choose to report decides
whether fine-tuning ``looks like'' a win}. We then show an early, cheap \emph{remedy}: instead
of tuning harder, keep the untuned base in the decision --- averaging the base and the tuned
guard's scores (a two-model ``ensemble'') keeps almost all the in-domain gain \emph{and}
recovers much of the lost generalization, most where fine-tuning hurt most. Every concept
needed to read this (guard scoring, LoRA, average precision, the bootstrap, calibration,
ensembling) is taught inline. The numbers are a retrospective ``legacy'' run and are reported
as estimates with uncertainty ranges, not as proven claims; a clean re-run would be needed to
confirm them. This is a plain-language companion to the formal paper, which remains the source
of truth wherever this edition simplifies.
\end{abstract}

\setcounter{tocdepth}{1}
\tableofcontents

\clearpage
% ============================================================
\section{Introduction}
% ============================================================

A \textbf{prompt-safety guard} reads a user's prompt and decides whether it is okay
(\textbf{safe}) or an attack or harmful request (\textbf{unsafe}). ``Unsafe'' here means one
of three things: \emph{harmful content} (asking for something dangerous), a \emph{jailbreak}
(trying to trick the model into ignoring its rules), or a \emph{prompt injection} (hiding
instructions in the input to hijack the system). Guards are attractive because they are
small and cheap to run in front of a bigger model.

The common recipe is: take a general chat model, fine-tune it on your own labeled prompts,
and report its score on a benchmark. The trouble is that a benchmark score, on its own,
hides the quantity a practitioner actually needs: \emph{what changed relative to the same
model before fine-tuning, and does that change generalize?} If you only compare your tuned
guard against \emph{other} models, a gain on the datasets you trained on can masquerade as a
broadly smarter guard.

\begin{background}{what a guard is and how we score it}
\small
The guard is just a language model asked to answer with one word: \code{safe} or
\code{unsafe}. Instead of letting it write a sentence, we do \textbf{one forward pass} and
look only at the two candidate next-words. For each, the model produces a raw preference
number (a \textbf{logit}). The score is the gap between them,
\[
s(x) = z_{\text{unsafe}} - z_{\text{safe}},
\]
positive meaning ``leans unsafe.'' A \textbf{softmax} over just those two turns the pair into
a probability that adds up to 1. \emph{Mini-example:} if $z_{\text{unsafe}}=2.0$ and
$z_{\text{safe}}=1.0$, then $s=1.0$ and
$P(\text{unsafe}) = e^{2.0}/(e^{2.0}+e^{1.0}) \approx 0.73$. Ranking prompts by this score is
all the guard has to do.
\end{background}

\paragraph{The question.} For a fixed panel of four checkpoints and one frozen LoRA recipe,
does fine-tuning change a guard's discrimination (how well it separates unsafe prompts from
safe ones) \emph{differently} on datasets represented during training versus datasets held out
from training?

\paragraph{Contributions.} (1) A controlled, same-model comparison: every fine-tuned guard is
measured against \emph{its own} untuned base on identical prompts, metrics, and scoring.
(2) A separation of two effects that are usually blended into one leaderboard number: gains
on \emph{represented} sources versus \emph{transfer} to held-out datasets, with per-model and
per-benchmark detail. (3) A fully auditable release: hashes, seed-level results, the scoring
table, and generated tables/figures.

\paragraph{Scope.} English input-prompt classification only; a binary safe/unsafe decision;
small models (1.5--4B); one LoRA recipe; three represented datasets; four transfer datasets;
two single-class stress sets; four fixed checkpoints. It is \emph{not} a leaderboard, a
fine-tuning-objective comparison, an ensembling or fairness study, or a domain-compliance
case study.

% ============================================================
\section{Background you'll need}
% ============================================================
This section teaches the five ideas the results rest on. If you already know average
precision, the bootstrap, and calibration, skip to Section~\ref{sec:design}.

\subsection{LoRA supervised fine-tuning (SFT)}
\begin{background}{LoRA-SFT}
\small
LoRA fine-tuning does \textbf{not} create a new model. The original model is \textbf{frozen}
(no weight changes). Instead you bolt on a small set of extra trainable weights --- an
\textbf{adapter} --- plugged into specific matrices inside the model. Only the adapter learns,
so it is cheap: you train a small fraction of extra parameters (millions, not billions) and
ship a small file that snaps onto the frozen base. ``Supervised'' means training on labeled
examples (prompt $+$ correct verdict word). \textbf{Completion-only loss on the verdict
token} means the model is graded \emph{only} on producing the single answer token, not on
echoing the prompt --- so all the learning pressure lands on that one decision.
\end{background}

\subsection{Average precision (AP), and why not accuracy}
\label{subsec:ap}
\begin{background}{average precision}
\small
The guard outputs a \emph{score} (the logit difference), and a good guard gives
\emph{higher} scores to truly-unsafe prompts. Two everyday quantities describe how it does:
\textbf{precision} $=$ of the prompts flagged so far, the fraction that are truly unsafe; and
\textbf{recall} $=$ of all the unsafe prompts, the fraction we have flagged.
\textbf{Average Precision (AP)} rolls these up with \textbf{no threshold to pick}: it is the
area under the precision--recall curve (0--1, higher better) --- in words, how well the guard
\emph{ranks} unsafe prompts above safe ones. Why not plain \textbf{accuracy}? Accuracy needs a
fixed cutoff and can be a trap --- if 95\% of prompts are safe, a guard that labels
\emph{everything} safe scores 95\% accuracy while catching zero attacks. AP does not fall for
that. \emph{Mini-example:} two unsafe and three safe prompts, ranked
\{unsafe, safe, unsafe, safe, safe\}. Precision at the first unsafe $=1/1=1.0$ (1 unsafe among
the top 1); at the second $=2/3\approx0.67$ (2 unsafe among the top 3);
$\text{AP}=(1.0+2/3)/2\approx\mathbf{0.83}$ (a perfect ranking scores 1.0).
\emph{Tie-aware} means two prompts with the same score get fair, not lucky, credit.
\end{background}
\noindent Throughout, the headline metric is \textbf{benchmark-macro AP}: average the AP
across the datasets in a group so each dataset counts equally (``macro''), then average across
the four models, and for fine-tuned guards across the five seeds too.

\subsection{Represented sources vs.\ held-out transfer}
\label{subsec:regimes}
\begin{background}{the seen-vs-unseen distinction (the crux)}
\small
\textbf{Represented} sources are fresh, held-out rows from the \emph{same} datasets used in
training (same style and labeling habits, different examples). Like taking this year's exam
on a topic you crammed from last year's --- you do well, partly because you learned the test's
patterns. \textbf{Transfer} is a \emph{new-material exam}: four entirely different datasets
never used in training. Reporting a single mixed number hides the difference between ``got
better at \emph{these} datasets'' and ``got broadly better'' --- exactly the confusion this
paper is about. \emph{(Honest caveat: the transfer sets were looked at during development, so
they are ``dataset-held-out,'' not a sealed surprise.)}
\end{background}

\subsection{The bootstrap and confidence intervals; ``descriptive, not proven''}
\label{subsec:boot}
\begin{background}{bootstrap confidence intervals}
\small
Each headline number is one estimate from the few models and seeds we ran; a different draw
could land higher or lower. To see how much it would \emph{wobble}, we \textbf{resample our own
results} (with replacement) 10{,}000 times and recompute the average each time (a
\textbf{paired hierarchical bootstrap}: \emph{paired} = base vs.\ fine-tuned, to measure the
\emph{change}; \emph{hierarchical} = resample at two levels, which seeds count and --- because
the eval sets contain many near-duplicate prompts --- how much each group of duplicates
counts; the four models are held fixed). The middle 95\% of those recomputed averages is the
\textbf{two-sided 95\% confidence interval (CI)} --- a plausible range for the true change. A
one-sided \textbf{LCB} says ``at least this much''; a \textbf{UCB} says ``no more than this.''
\textbf{Why no ``statistically significant''?} These scores are a \emph{legacy} artifact
computed before the analysis was locked; declaring significance on already-peeked data is
misleading. So the study stays \textbf{descriptive}: it reports ranges, makes no pass/fail
verdict, and treats directions as \emph{suggestive} pending a clean, pre-registered re-run.
\end{background}

\subsection{Calibration and the operating point}
\label{subsec:calib}
\begin{background}{calibration, thresholds, TPR and FPR}
\small
AP is threshold-free, but to \emph{deploy} a guard you must draw a line. First,
\textbf{calibration (temperature scaling)}: divide the score by one tuned number so the stated
probabilities match reality (this does not reorder prompts). Second, a \textbf{threshold}:
here a \emph{conservative cap} chosen so at most ${\sim}5\%$ of genuinely safe prompts are
flagged (the target false-positive rate); because it is conservative, the realized rate can
land well below 5\%. Then you \emph{measure} on test data: \textbf{TPR} (recall) $=$ share of
unsafe prompts caught (higher better); \textbf{FPR} $=$ share of safe prompts wrongly flagged
(lower better). \emph{Mini-example (100 prompts, 80 safe / 20 unsafe):} a cutoff that flags
$\sim$5\% of safe prompts $=$ 4 of 80 (FPR $=5\%$); at that cutoff it catches 15 of 20 unsafe
(TPR $=75\%$).
\end{background}

% ============================================================
\section{Study design}
\label{sec:design}
% ============================================================

\subsection{The panel and the frozen recipe}
Four instruction models are each adapted with five training seeds (42--46), giving 20
fine-tuned guards; the four untuned bases are scored once and reused ($4+20=24$ ``bundles''
in total, and $79{,}392$ scored rows). Every cell uses the identical recipe in
Table~\ref{tab:panel}. Five seeds are used because a single fine-tune is slightly random;
running five shows how much the result wobbles from initialization and ordering alone (the
training rows and their order are fixed).

\begin{table}[t]
\centering\small
\caption{The fixed four-checkpoint panel and the one LoRA-SFT recipe used for every cell.}
\label{tab:panel}
\begin{tabular}{lll}
\toprule
Checkpoint & Params & Role \\
\midrule
Qwen2.5-1.5B-Instruct & 1.5B & base $\to$ 5 SFT seeds \\
SmolLM2-1.7B-Instruct & 1.7B & base $\to$ 5 SFT seeds \\
SmolLM3-3B            & 3B   & base $\to$ 5 SFT seeds \\
Qwen3-4B             & 4B   & base $\to$ 5 SFT seeds \\
\midrule
\multicolumn{3}{l}{\emph{Recipe:} LoRA $r{=}32$, $\alpha{=}64$, dropout $0.05$ on
q,k,v,o,gate,up,down; completion-only loss;}\\
\multicolumn{3}{l}{\phantom{\emph{Recipe:}} $1{,}200$ training rows; 300 steps; learning rate
$2\times10^{-4}$ cosine; effective batch 4; max length 1024.}\\
\bottomrule
\end{tabular}
\end{table}

\subsection{Training data, label mapping, and decontamination}
Training reads one frozen manifest of \textbf{1{,}200 rows}: 400 from each of three datasets
--- ToxicChat, Prompt-Injections, and Jailbreak-Classification --- with 200 safe and 200
unsafe per dataset, selected by a fixed hash rank so the row set is identical across every
model and seed. Each dataset's native label is mapped to safe/unsafe by a fixed rule (e.g.\
ToxicChat \code{toxicity=1}$\to$unsafe; Prompt-Injections \code{label=1}$\to$unsafe;
Jailbreak-Classification \code{type} starting ``jailbreak''$\to$unsafe). These are
\emph{different} native policies, not one universal definition of ``unsafe.''

To keep the transfer test honest, the builder removes exact and near-duplicate overlaps
between training rows and any evaluation row (grouping paraphrases into ``families'' via
MinHash and refusing to let a family straddle the train/eval boundary). This is why a
``transfer'' gain or loss cannot be explained away as the model having memorized a training
row that reappears in the test set.

\begin{table}[t]
\centering\small
\caption{Datasets and their roles. Represented sources appear (as different rows) in both
training and the represented test; transfer and stress sets are test-only.}
\label{tab:data}
\begin{tabular}{lll}
\toprule
Group & Datasets & Use \\
\midrule
Represented & ToxicChat, Prompt-Injections, Jailbreak-Classification & train $+$ represented test \\
Transfer & JailbreakBench, XSTest, WildGuardTest, WildJailbreak & test only (never trained on) \\
Stress & OR-Bench (all benign), HarmBench (all harmful) & one-sided diagnostics only \\
\bottomrule
\end{tabular}
\end{table}

\subsection{The three test groups}
\begin{itemize}
\item \textbf{Represented} $=$ held-out rows from the three training datasets --- measures
discrimination on \emph{sources seen in training}.
\item \textbf{Transfer} $=$ four datasets withheld from training --- measures whether the skill
\emph{generalizes} (Section~\ref{subsec:regimes}).
\item \textbf{Stress} $=$ two single-class sets. OR-Bench is all benign (measures false
alarms); HarmBench is all harmful (measures catch rate). Because each has only one class,
\textbf{no AP is computed on them} --- they are diagnostics.
\end{itemize}

\subsection{What we compute}
For a group $R$, checkpoint $b$, and seed $r$, the metric is the benchmark-macro AP
(Section~\ref{subsec:ap}). The change from fine-tuning is the paired difference
\[
\Delta_R(b,r) = \text{macroAP}_R(\text{SFT}_{b,r}) - \text{macroAP}_R(\text{base}_b),
\]
and the panel aggregate averages the four checkpoints (averaging seeds within each):
\[
\bar\Delta_R = \frac{1}{4}\sum_b\Big[\tfrac{1}{5}\sum_r \text{macroAP}_R(\text{SFT}_{b,r}) -
\text{macroAP}_R(\text{base}_b)\Big].
\]
The primary result is the \emph{pair}
$(\bar\Delta_{\text{represented}}, \bar\Delta_{\text{transfer}})$ --- we keep both numbers
rather than collapsing them into one ``specialization score.'' Uncertainty is the paired
hierarchical bootstrap of Section~\ref{subsec:boot}; the analysis mode is
\emph{descriptive} (no significance test).

% ============================================================
\section{Results}
% ============================================================

\subsection{Represented sources: a big, reliable win}
On the datasets the models trained on, every fine-tuned guard reaches ${\approx}0.98$
macro-AP regardless of where its base started (Table~\ref{tab:primary}, left). The panel
aggregate change is $\bar\Delta_{\text{represented}} = \mathbf{+0.3327}$ (95\% CI
$[+0.2718, +0.3793]$). Notice that the \emph{weakest} base gains the most (SmolLM2:
$0.447\!\to\!0.981$, $+0.534$) and the \emph{strongest} gains the least (Qwen3-4B:
$0.878\!\to\!0.983$, $+0.104$) --- largely a ceiling effect: everyone ends near 0.98, so
whoever started lowest had the most room to climb. Per-dataset gains: ToxicChat $+0.19$,
Prompt-Injections $+0.38$, Jailbreak-Classification $+0.44$ (Table~\ref{tab:perbench}).

\subsection{Transfer: a small average loss that hides a split}
On the four never-trained-on datasets, the same fine-tuning is a small average loss:
$\bar\Delta_{\text{transfer}} = \mathbf{-0.0503}$ (95\% CI $[-0.0760, -0.0250]$). But the
per-checkpoint column of Table~\ref{tab:primary} shows the average is misleading: one model
\emph{improves}, one is uncertain, and two clearly drop. The losses concentrate on the
jailbreak-style sets (WildJailbreak $-0.079$, JailbreakBench $-0.073$), while the
over-refusal set barely moves (XSTest $-0.010$); see Table~\ref{tab:perbench}.

\begin{table}[t]
\centering\small
\caption{Primary result: base and mean-SFT benchmark-macro AP per checkpoint, with the paired
change $\Delta$ and its two-sided 95\% bootstrap interval. Positive $\Delta$ $=$ fine-tuning
helped. The last row is the panel aggregate.}
\label{tab:primary}
\setlength{\tabcolsep}{4pt}
\begin{tabular}{l rr r@{\ }l rr r@{\ }l}
\toprule
& \multicolumn{3}{c}{Represented} && \multicolumn{3}{c}{Transfer} &\\
\cmidrule(lr){2-4}\cmidrule(lr){6-8}
Checkpoint & base & SFT & $\Delta$ & [95\% CI] & base & SFT & $\Delta$ & [95\% CI]\\
\midrule
Qwen2.5-1.5B & 0.622 & 0.988 & $+0.366$ & [0.283, 0.426] & 0.822 & 0.791 & $-0.030$ & [$-0.078$, $+0.018$]\\
SmolLM2-1.7B & 0.447 & 0.981 & $+0.534$ & [0.461, 0.580] & 0.787 & 0.838 & $+0.051$ & [0.018, 0.082]\\
SmolLM3-3B  & 0.655 & 0.982 & $+0.327$ & [0.253, 0.385] & 0.914 & 0.794 & $-0.120$ & [$-0.156$, $-0.083$]\\
Qwen3-4B   & 0.878 & 0.983 & $+0.104$ & [0.060, 0.154] & 0.945 & 0.843 & $-0.102$ & [$-0.129$, $-0.075$]\\
\midrule
\textbf{Panel aggregate} & -- & -- & $+0.333$ & [0.272, 0.379] & -- & -- & $-0.050$ & [$-0.076$, $-0.025$]\\
\bottomrule
\end{tabular}
\end{table}

\begin{table}[t]
\centering\small
\caption{Per-dataset paired AP change (base $\to$ SFT). Represented gains are large and
uniform; transfer losses concentrate on jailbreak-style sets.}
\label{tab:perbench}
\begin{tabular}{llr}
\toprule
Group & Dataset & $\Delta$ AP \\
\midrule
Represented & ToxicChat & $+0.187$ \\
Represented & Prompt-Injections & $+0.375$ \\
Represented & Jailbreak-Classification & $+0.436$ \\
\midrule
Transfer & JailbreakBench & $-0.073$ \\
Transfer & XSTest & $-0.010$ \\
Transfer & WildGuardTest & $-0.039$ \\
Transfer & WildJailbreak & $-0.079$ \\
\bottomrule
\end{tabular}
\end{table}

\subsection{The key nuance: who fine-tuning helps vs.\ hurts}
Read the ``SFT'' columns of Table~\ref{tab:primary}: on transfer, fine-tuning squeezes every
model into a narrow $\approx0.79$--$0.84$ band, even though the bases ranged from 0.79 to 0.94.
Because that band acts like a fixed destination, fine-tuning has room to \emph{pull up} a weak
base and little room to help a strong one --- so it \emph{drags strong bases down}.

\begin{takeaway}{Why the $-0.05$ average is misleading}
\small
Imagine fine-tuning always lands new-data skill at ${\sim}0.80$. A weak model at 0.70 gains
$+0.10$; a strong one at 0.95 loses $-0.15$. Average them: $(+0.10-0.15)/2 = -0.025$ --- a
tiny dip that \emph{erases both the real help and the real harm}. In the data, the weakest base
(SmolLM2) is the only transfer gainer ($+0.051$); the strongest (SmolLM3 $-0.120$, Qwen3-4B
$-0.102$) lose the most. \emph{(These per-model directions are descriptive point estimates from
a legacy run --- suggestive, not proven.)}
\end{takeaway}

\subsection{The specialization plane}
Figure~\ref{fig:plane} puts both changes on one chart: each dot is one (checkpoint, seed);
right means represented AP went up, up means transfer AP went up. \textbf{14 of the 20 dots}
land in the lower-right ``specialization'' quadrant (better on familiar data, worse on new);
the other 6 are in ``uniform gain'' (all five SmolLM2 seeds plus one Qwen2.5 seed), and none
are in a loss quadrant. The dominant story is specialization, and the exceptions are exactly
the weak base.

\begin{figure}[t]
\centering
\includegraphics[width=0.82\linewidth]{specialization_plane.pdf}
\caption{The specialization plane. Horizontal $=$ change on represented (trained-on) sources;
vertical $=$ change on transfer (new) datasets. Each dot is one (checkpoint, seed); the
$\times$ is the panel mean. All 20 dots move right (represented up); 14 also move down
(transfer down).}
\label{fig:plane}
\end{figure}

\subsection{At a deployed threshold (operating point)}
AP is threshold-free; Table~\ref{tab:op} shows what happens once you pick a cutoff at a
${\sim}5\%$ target false-positive rate (Section~\ref{subsec:calib}). On represented data the
guard becomes far more useful --- recall jumps $13.4\%\!\to\!76.6\%$ at only ${\sim}1\%$ false
alarms. On transfer data recall barely moves ($52.6\%\!\to\!55.4\%$) while false alarms get
\emph{worse}: benchmark-macro FPR $8.3\%\!\to\!13.7\%$, pooled FPR $4.4\%\!\to\!14.6\%$
(\emph{pooled} $=$ every test row combined into one set, so larger datasets dominate;
\emph{macro} $=$ each dataset weighted equally --- which is why the base's pooled $4.4\%$ and
macro $8.3\%$ differ).

\begin{takeaway}{Why the transfer FPR (13.7\%) exceeds the 5\% cap}
\small
The ${\sim}5\%$ cap is only guaranteed on the \emph{calibration} data. The cutoff is frozen
there and never re-tuned, so on a new distribution it no longer holds and the realized false
alarm rate sails past 5\% (the base is already at 8.3\%). Losing threshold control on new data
is itself part of the finding.
\end{takeaway}

\subsection{Stress sets: false alarms and missed attacks}
The two single-class diagnostics (Table~\ref{tab:op}, lower panel) sharpen the picture.
Fine-tuning slightly \emph{reduces} over-blocking of benign prompts (OR-Bench FPR
$12.7\%\!\to\!10.0\%$) --- a small plus. But on hard held-out harmful prompts it catches
\emph{fewer}: HarmBench recall falls $78.4\%\!\to\!57.5\%$, a \textbf{20.9-point drop}. That is
the sharpest downside in the study, and a safety-relevant one: the specialized guard misses
about one in five more real attacks it would otherwise have caught.

\begin{table}[t]
\centering\small
\caption{Deployed-threshold behavior (calibration-targeted ${\sim}5\%$ FPR) and single-class
stress diagnostics. TPR $=$ recall; FPR $=$ false-alarm rate; $N$ $=$ rows in the relevant class.}
\label{tab:op}
\begin{tabular}{ll rr r r}
\toprule
Group & Measure & base & SFT & change & $N$ \\
\midrule
Represented & TPR (recall) & 13.4\% & 76.6\% & $+63.2$ & 312 \\
Represented & FPR (macro)  & 1.7\%  & 1.0\%  & $-0.7$  & 364 \\
Transfer    & TPR (recall) & 52.6\% & 55.4\% & $+2.8$  & 790 \\
Transfer    & FPR (macro)  & 8.3\%  & 13.7\% & $+5.4$  & 790 \\
Transfer    & FPR (pooled) & 4.4\%  & 14.6\% & $+10.2$ & 790 \\
\midrule
Stress: OR-Bench & benign FPR & 12.7\% & 10.0\% & $-2.7$ & 400 \\
Stress: HarmBench & recall    & 78.4\% & 57.5\% & $-20.9$ & 200 \\
\bottomrule
\end{tabular}
\end{table}

% ============================================================
\section{A fix: compose, don't tune}
\label{sec:fix}
% ============================================================
The results so far pose a problem: fine-tuning wins big on familiar data but gives up ground
on new data and misses more hard attacks (\Cref{tab:op}). The obvious reaction is to fine-tune
\emph{better}. Here is an early look at a cheaper fix (from a planned follow-up, ``Compose,
Don't Tune''): instead of tuning harder, \textbf{keep the original untuned model in the decision}.

\begin{background}{combining two guards (an ``ensemble'')}
\small
Run \emph{two} guards on each prompt --- the untuned base and the fine-tuned adapter --- and
average their unsafe-probabilities into one score. (Using more than one model and combining
their outputs is called an \textbf{ensemble}.) Think of it as a \emph{second opinion}: the
fine-tuned specialist is best on prompts like the ones it trained on; the untuned generalist
holds up better on new kinds of prompts; averaging keeps much of both. It is cheap --- the two
guards share the same backbone with the adapter switched on or off, so it is two quick passes
and \textbf{no extra training}. \emph{Mini-example:} on a new prompt the base scores $0.30$
unsafe and the tuned guard scores $0.90$; the ensemble scores $(0.30+0.90)/2 = 0.60$.
\end{background}

Running this on the same data as before (\Cref{tab:compose}): the ensemble keeps almost all of
the tuned guard's gain on familiar data ($0.99 \to 0.97$) \emph{and} does clearly better than
the tuned guard on new data ($0.84 \to 0.90$), pulling generalization back up toward the untuned
base --- without any extra fine-tuning.

\begin{table}[t]
\centering\small
\caption{Composing the base with the tuned guard. Benchmark-macro AP (0--1, higher better),
averaged over the four checkpoints. The ensemble keeps the tuned guard's familiar-data gain and
recovers most of its lost new-data score.}
\label{tab:compose}
\begin{tabular}{lrr}
\toprule
Guard & Familiar (represented) & New (transfer) \\
\midrule
Untuned base & 0.65 & 0.87 \\
Tuned (SFT) & \textbf{0.99} & 0.84 \\
Composed (base $+$ tuned, averaged) & 0.97 & \textbf{0.90} \\
\bottomrule
\end{tabular}
\end{table}

\begin{takeaway}{What composing does --- and doesn't do}
\small
Composing \textbf{reliably beats the tuned guard} on new data --- for every model in the panel.
It does \textbf{not} beat a \emph{strong untuned base}: for the strongest model (Qwen3-4B) the
ensemble is a touch \emph{below} the base on new data (0.93 vs 0.95), and for SmolLM3 it only
ties. So composing is best understood as \textbf{protecting against the damage fine-tuning did},
not a free lunch that beats everything. Tellingly, it helps most exactly where fine-tuning hurt
most --- the strong bases (\Cref{tab:composeperm}).
\end{takeaway}

Read \Cref{tab:composeperm} across each row: the tuned column is always at or below the untuned
base on new data (that is the specialization cost from \Cref{tab:primary}); the composed column
is always at or above the tuned one, and lands back near the base --- most dramatically for
SmolLM3 (0.82 tuned, back to 0.91) and Qwen3-4B (0.87 tuned, back to 0.93).

\begin{table}[t]
\centering\small
\caption{Per-model new-data (transfer) macro-AP. Fine-tuning lowers the base (specialization);
composing recovers it. The gain from composing is largest for the strong bases fine-tuning hurt
most.}
\label{tab:composeperm}
\begin{tabular}{lrrr}
\toprule
Model (new-data strength of base) & Untuned base & Tuned & Composed \\
\midrule
SmolLM2-1.7B (weakest) & 0.79 & 0.86 & 0.87 \\
Qwen2.5-1.5B & 0.82 & 0.82 & 0.87 \\
SmolLM3-3B & 0.91 & 0.82 & 0.91 \\
Qwen3-4B (strongest) & 0.95 & 0.87 & 0.93 \\
\bottomrule
\end{tabular}
\end{table}

\begin{background}{is this real, or just a trick of averaging?}
\small
A fair worry: does averaging two scores help \emph{automatically}, no matter what? Not here.
Averaging only helps when the two guards make \emph{different} mistakes --- the specialist and
the generalist are each right on prompts the other gets wrong, so combining them cancels errors;
if they agreed on everything, averaging would change nothing. A planned control tests exactly
this (break the two guards' complementarity and the ensemble's advantage should vanish),
confirming the gain is genuine \emph{teamwork}, not a free bump from averaging any two numbers.
\end{background}

\paragraph{What this fix does not solve (yet).} Three honest caveats. (1) These are the same
\emph{legacy} scores as the rest of the paper, so treat the composition numbers as an early
signal, not a confirmed result --- a clean re-run is needed. (2) Composing improves the
\emph{ranking} of prompts, not the \emph{calibration}: at a fixed yes/no cutoff the false-alarm
rate on new data is still too high, so you would still re-tune the threshold for new traffic.
(3) It costs two passes per prompt instead of one (still cheap, but not free).

% ============================================================
\section{What this means in practice}
% ============================================================
\begin{takeaway}{Practitioner takeaways}
\small
\begin{itemize}
\item \textbf{Fine-tune when your live traffic resembles your training sources.}
In-distribution, the win is large and reliable (recall $13\%\!\to\!77\%$).
\item \textbf{Expect erosion on novel attacks.} On new datasets, fine-tuning trades away
generalization --- more false alarms and a real drop in catching hard, unseen harm.
\item \textbf{Always test on datasets the model never trained on.} A single trained-on
benchmark will crown fine-tuning the winner and hide the transfer cost. Report \emph{both}
groups.
\item \textbf{Mind the base's starting strength.} A weak base has the most to gain; a strong
base may be better left \emph{untouched} for screening new/out-of-distribution traffic --- it
can already be the better guard there. \emph{(A suggestive 4-model pattern, not a proven law.)}
\item \textbf{Prefer composing over tuning harder.} If you do fine-tune, keep the untuned base
in the decision --- average the two guards' scores (\Cref{sec:fix}). On this panel that keeps
almost all the in-domain gain and recovers much of the lost generalization, for two cheap passes
and no extra training.
\end{itemize}
\end{takeaway}

% ============================================================
\section{Limitations (what not to over-read)}
% ============================================================
\begin{itemize}
\item \textbf{Legacy artifact.} The scores predate the locked pipeline, so results are
\emph{estimation-only}; a clean, pre-registered re-run is needed before any number is
``confirmatory.''
\item \textbf{Balanced test pools overstate real-world precision.} Test sets are ${\sim}50/50$
safe/unsafe; real traffic has far fewer unsafe prompts, so precision/AP look rosier here than
in production.
\item \textbf{Transfer sets were not truly sealed.} They were inspected during development, so
``transfer'' means dataset-held-out, not never-before-seen.
\item \textbf{Only four models, two families (Qwen, SmolLM), 1.5--4B.} Conclusions are about
\emph{this panel}; the base-competence pattern rests on four points.
\item \textbf{Mixed native policies.} The four transfer datasets encode different definitions
of ``unsafe'' (harm vs.\ refusal vs.\ jailbreak), so their macro-average blends distinct
policies.
\item \textbf{Prompt-only, English.} No response, tool-call, or multi-turn moderation.
\item \textbf{The ``compose'' fix (\Cref{sec:fix}) is preliminary.} It runs on the same legacy
scores, improves ranking but not calibration at a fixed cutoff, and costs two passes per prompt;
a clean re-run is needed to confirm it.
\end{itemize}

% ============================================================
\section{Conclusion}
% ============================================================
Fine-tuning \emph{specializes} a small safety guard: large, reliable gains on data resembling
its training sources (represented macro-AP $+0.33$, everyone to ${\approx}0.98$), a small
average loss on genuinely new datasets ($-0.05$), and a real drop in catching hard unseen harm
(HarmBench recall $-20.9$ points). The average transfer number hides a split --- fine-tuning
helps weak bases and hurts strong ones --- and 14 of 20 fine-tunes land in the
``better-on-familiar, worse-on-new'' quadrant. So the benchmark you choose to report decides
whether fine-tuning looks like a win. Report \emph{both} groups with uncertainty ranges, and
treat these specific numbers as \emph{suggestive pending a clean re-run}. And there is an early,
cheap way out of the trade-off: rather than tuning harder, \emph{compose} the tuned guard with
its untuned base (\Cref{sec:fix}) --- on this panel that keeps the in-domain gains while
recovering much of the lost generalization.

% ============================================================
\appendix
\section{Per-seed changes}
% ============================================================
Table~\ref{tab:seeds} lists the paired change for every (checkpoint, seed), the raw material
behind the panel aggregates and Figure~\ref{fig:plane}. The seed-to-seed spread is small on
represented sources and larger on transfer --- another reason the study reports ranges, not
single points.

\begin{table}[h]
\centering\small
\caption{Represented and transfer AP change ($\Delta$) for each of the 20 fine-tunes.}
\label{tab:seeds}
\begin{tabular}{l rrrrr}
\toprule
& \multicolumn{5}{c}{seed} \\
\cmidrule(lr){2-6}
Checkpoint (group) & 42 & 43 & 44 & 45 & 46 \\
\midrule
Qwen2.5-1.5B (rep.)  & $+0.366$ & $+0.368$ & $+0.365$ & $+0.362$ & $+0.367$ \\
Qwen2.5-1.5B (tr.)   & $-0.079$ & $+0.037$ & $-0.046$ & $-0.037$ & $-0.026$ \\
SmolLM2-1.7B (rep.)  & $+0.536$ & $+0.533$ & $+0.530$ & $+0.533$ & $+0.537$ \\
SmolLM2-1.7B (tr.)   & $+0.073$ & $+0.075$ & $+0.034$ & $+0.038$ & $+0.036$ \\
SmolLM3-3B (rep.)   & $+0.328$ & $+0.328$ & $+0.327$ & $+0.327$ & $+0.325$ \\
SmolLM3-3B (tr.)    & $-0.141$ & $-0.075$ & $-0.133$ & $-0.146$ & $-0.106$ \\
Qwen3-4B (rep.)    & $+0.104$ & $+0.108$ & $+0.105$ & $+0.102$ & $+0.102$ \\
Qwen3-4B (tr.)     & $-0.114$ & $-0.086$ & $-0.125$ & $-0.087$ & $-0.100$ \\
\bottomrule
\end{tabular}
\end{table}

\section{Provenance and reproducibility}
The scored data table stores per-row \emph{content hashes and model logits only --- never raw
prompt text} --- so third-party benchmark text is not redistributed, yet every metric is
recomputable. Metrics come from one shared, tie-aware module (no ad-hoc AP loop). Source
datasets are pinned by revision; near-duplicate ``families'' are detected with a fixed MinHash
so the split cannot change with the environment. Full identifiers, hashes, licenses, and
seed-level outputs are in \code{artifacts/paper\_a\_sft/}.

\section*{Sources and further reading}
\small
Guard datasets and benchmarks referenced above: ToxicChat, Prompt-Injections
(\code{deepset}), Jailbreak-Classification, JailbreakBench, XSTest, WildGuardMix /
WildGuardTest, WildJailbreak, OR-Bench, and HarmBench. Methods: LoRA (low-rank adaptation) for
parameter-efficient fine-tuning; temperature scaling for calibration; average precision as a
threshold-free ranking metric. Related guards include Llama Guard, ShieldGemma, Granite
Guardian, and WildGuard. For full citations, the formal statistical protocol, and the exact
lock/audit contract, see the formal paper under \code{paper-a/} and its reference list.

\end{document}
